% GENERATED FILE. Edit canonical lessons, then run npm run generate:books.
% canonical-source-hash: fnv1a64:e652541fdceb08b8
% canonical-lessons: FA-C496-maghz, FA-C496-asab, FA-C496-rag, FA-C496-pelk, FA-C496-lase

\chapter{Brain, Nerves and Eyes}
\label{ch:fa-brain-nerves-and-eyes}
\hlchaptermodality{496}

\begin{chapteropening}
\textbf{By the end of this chapter:} \emph{I can talk about the brain, nerves and eyes.}

Complete the last lesson of chapter 496: I can talk about the brain, nerves and eyes.
\end{chapteropening}

\section[\texorpdfstring{\fa{مغز} (maghz)}{مغز (maghz)}]{\fa{مغز} (maghz) — the brain}

\label{lesson:FA-C496-maghz}



\begin{quote}
Before the new one: say the Persian for a kidney, then the Persian for the stomach.
\end{quote}

\subsection*{You'll want to know: \fa{مغز}}
\textbf{\fa{مغز}} (\emph{maghz}) — \textquotedblleft{}the brain\textquotedblright{}.

\begin{quote}
\fa{مغز} \fa{در} \fa{خواب} \fa{هم} \fa{کار} \fa{میکند}.
\end{quote}

\begin{quote}
\emph{maghz dar khâb ham kâr mikonad.}
\end{quote}

\begin{quote}
The brain works even during sleep.
\end{quote}

Also the kernel of a nut: \emph{maghz-e gerdu} 'a walnut kernel'.

\begin{cousinweb}
\emph{maghz} is inherited: EtymDB traces it to Middle Persian, where the word meant 'brain, marrow' (EtymDB).
\end{cousinweb}

\begin{grammarlens}[title={What you've built}]
One more word for telling, explaining and describing at B1.
\end{grammarlens}

\subsection*{Your turn}
\begin{itemize}
\raggedright
  \item \emph{Say it:} \emph{maghz}
  \item \emph{Say it:} \emph{maghz}, once more
  \item \emph{Say it:} a sentence of your own with \emph{maghz}
  \item \emph{Recall:} say the Persian for a kidney, then the Persian for the stomach, then say \emph{maghz} again
\end{itemize}

\begin{tcolorbox}[breakable,colback=teal!4,colframe=teal!35!black,title={Before you move on}]
What does \fa{مغز} mean? (\textquotedblleft{}The brain\textquotedblright{}.) Say it once more.
\end{tcolorbox}
\section[\texorpdfstring{\fa{اعصاب} (a'sâb)}{اعصاب (a'sâb)}]{\fa{اعصاب} (a'sâb) — nerves}

\label{lesson:FA-C496-asab}



\begin{quote}
Before the new one: say the Persian for the stomach, then the Persian for the brain.
\end{quote}

\subsection*{You'll want to know: \fa{اعصاب}}
\textbf{\fa{اعصاب}} (\emph{a'sâb}) — \textquotedblleft{}nerves\textquotedblright{}.

\begin{quote}
\fa{صدای} \fa{بلند} \fa{روی} \fa{اعصاب} \fa{من} \fa{است}.
\end{quote}

\begin{quote}
\emph{sedâ-ye boland ru-ye a'sâb-e man ast.}
\end{quote}

\begin{quote}
Loud noise gets on my nerves.
\end{quote}

A plural noun. \emph{ru-ye a'sâb budan} — to get on someone's nerves.

\begin{cousinweb}
\emph{a'sâb} is a loan, usually traced to Arabic \emph{aʿṣāb}, the plural of \emph{ʿaṣab}, 'nerve, sinew'. The same \emph{'asab} is inside \emph{'asabâni}. EtymDB has no entry for it.
\end{cousinweb}

\begin{grammarlens}[title={What you've built}]
One more word for telling, explaining and describing at B1.
\end{grammarlens}

\subsection*{Your turn}
\begin{itemize}
\raggedright
  \item \emph{Say it:} \emph{a'sâb}
  \item \emph{Say it:} \emph{a'sâb}, once more
  \item \emph{Say it:} a sentence of your own with \emph{a'sâb}
  \item \emph{Recall:} say the Persian for the stomach, then the Persian for the brain, then say \emph{a'sâb} again
\end{itemize}

\begin{tcolorbox}[breakable,colback=teal!4,colframe=teal!35!black,title={Before you move on}]
What does \fa{اعصاب} mean? (\textquotedblleft{}Nerves\textquotedblright{}.) Say it once more.
\end{tcolorbox}
\section[\texorpdfstring{\fa{رگ} (rag)}{رگ (rag)}]{\fa{رگ} (rag) — a vein, a blood vessel}

\label{lesson:FA-C496-rag}



\begin{quote}
Before the new one: say the Persian for the brain, then the Persian for nerves.
\end{quote}

\subsection*{You'll want to know: \fa{رگ}}
\textbf{\fa{رگ}} (\emph{rag}) — \textquotedblleft{}a vein, a blood vessel\textquotedblright{}.

\begin{quote}
\fa{پرستار} \fa{رگ} \fa{دستش} \fa{را} \fa{پیدا} \fa{نکرد}.
\end{quote}

\begin{quote}
\emph{parastâr rag-e dastash râ peydâ nakard.}
\end{quote}

\begin{quote}
The nurse couldn't find a vein in his arm.
\end{quote}

Plural \emph{rag-hâ}. Also a streak or vein in stone or wood.

\begin{cousinweb}
\emph{rag} is inherited: EtymDB traces it to Middle Persian (EtymDB).
\end{cousinweb}

\begin{grammarlens}[title={What you've built}]
One more word for telling, explaining and describing at B1.
\end{grammarlens}

\subsection*{Your turn}
\begin{itemize}
\raggedright
  \item \emph{Say it:} \emph{rag}
  \item \emph{Say it:} \emph{rag}, once more
  \item \emph{Say it:} a sentence of your own with \emph{rag}
  \item \emph{Recall:} say the Persian for the brain, then the Persian for nerves, then say \emph{rag} again
\end{itemize}

\begin{tcolorbox}[breakable,colback=teal!4,colframe=teal!35!black,title={Before you move on}]
What does \fa{رگ} mean? (\textquotedblleft{}A vein, a blood vessel\textquotedblright{}.) Say it once more.
\end{tcolorbox}
\section[\texorpdfstring{\fa{پلک} (pelk)}{پلک (pelk)}]{\fa{پلک} (pelk) — an eyelid}

\label{lesson:FA-C496-pelk}



\begin{quote}
Before the new one: say the Persian for nerves, then the Persian for a vein.
\end{quote}

\subsection*{You'll want to know: \fa{پلک}}
\textbf{\fa{پلک}} (\emph{pelk}) — \textquotedblleft{}an eyelid\textquotedblright{}.

\begin{quote}
\fa{او} \fa{بدون} \fa{پلک} \fa{زدن} \fa{به} \fa{من} \fa{نگاه} \fa{کرد}.
\end{quote}

\begin{quote}
\emph{u bedun-e pelk zadan be man negâh kard.}
\end{quote}

\begin{quote}
She looked at me without blinking.
\end{quote}

\emph{pelk zadan} — to blink.

\begin{cousinweb}
\emph{pelk} is usually taken as a native Persian word, but EtymDB lists it without any source, and its origin is not settled.
\end{cousinweb}

\begin{grammarlens}[title={What you've built}]
One more word for telling, explaining and describing at B1.
\end{grammarlens}

\subsection*{Your turn}
\begin{itemize}
\raggedright
  \item \emph{Say it:} \emph{pelk}
  \item \emph{Say it:} \emph{pelk}, once more
  \item \emph{Say it:} a sentence of your own with \emph{pelk}
  \item \emph{Recall:} say the Persian for nerves, then the Persian for a vein, then say \emph{pelk} again
\end{itemize}

\begin{tcolorbox}[breakable,colback=teal!4,colframe=teal!35!black,title={Before you move on}]
What does \fa{پلک} mean? (\textquotedblleft{}An eyelid\textquotedblright{}.) Say it once more.
\end{tcolorbox}
\section[\texorpdfstring{\fa{لثه} (lase)}{لثه (lase)}]{\fa{لثه} (lase) — the gum (in the mouth)}

\label{lesson:FA-C496-lase}



\begin{quote}
Before the new one: say the Persian for a vein, then the Persian for an eyelid.
\end{quote}

\subsection*{You'll want to know: \fa{لثه}}
\textbf{\fa{لثه}} (\emph{lase}) — \textquotedblleft{}the gum (in the mouth)\textquotedblright{}.

\begin{quote}
\fa{دندانپزشک} \fa{لثه} \fa{او} \fa{را} \fa{معاینه} \fa{کرد}.
\end{quote}

\begin{quote}
\emph{dandânpezeshk lase-ye u râ mo'âyene kard.}
\end{quote}

\begin{quote}
The dentist examined her gums.
\end{quote}

Plural \emph{lasehâ}.

\begin{cousinweb}
\emph{lase} is a loan, usually traced to the Arabic word for the gum of the mouth. EtymDB has no entry for it.
\end{cousinweb}

\begin{grammarlens}[title={What you've built}]
One more word for telling, explaining and describing at B1.
\end{grammarlens}

\subsection*{Your turn}
\begin{itemize}
\raggedright
  \item \emph{Say it:} \emph{lase}
  \item \emph{Say it:} \emph{lase}, once more
  \item \emph{Say it:} a sentence of your own with \emph{lase}
  \item \emph{Recall:} say the Persian for a vein, then the Persian for an eyelid, then say \emph{lase} again
\end{itemize}

\begin{tcolorbox}[breakable,colback=teal!4,colframe=teal!35!black,title={Before you move on}]
What does \fa{لثه} mean? (\textquotedblleft{}The gum (in the mouth)\textquotedblright{}.) Say it once more.
\end{tcolorbox}
